\chapter{Twierdzenie o zwartości i jego zastosowania}\label{ch:compactness}

\section{Dlaczego nieskończoną teorię można badać skończonymi fragmentami?}
Wyobraźmy sobie nieskończony zestaw wymagań wobec struktury. Nie potrafimy sprawdzić wszystkich naraz, ale każdy \emph{skończony} zestaw wymagań da się spełnić. Czy musi istnieć jedna struktura spełniająca je wszystkie? W dowolnej dziedzinie matematyki odpowiedź mogłaby być przecząca. Dla zdań klasycznej logiki pierwszego rzędu jest jednak twierdząca. To właśnie \emph{twierdzenie o zwartości}.

\begin{definicja}[Spełnialność i spełnialność skończona]\label{def:compactness-finite-satisfiability}
Niech $\mathcal L$ będzie językiem pierwszego rzędu, a $T$ zbiorem zdań tego języka. Mówimy, że $T$ jest \emph{spełnialny}, jeśli istnieje struktura $\mathfrak M$ taka, że $\mathfrak M\models T$. Mówimy, że $T$ jest \emph{skończenie spełnialny}, jeśli dla każdego skończonego podzbioru $T_0\subseteq T$ istnieje (być może zależny od $T_0$) model $\mathfrak M_{T_0}$ spełniający wszystkie zdania z $T_0$.
\end{definicja}

\begin{uwaga}[Model może zależeć od wybranych zdań]
Skończona spełnialność \emph{nie} zakłada, że wszystkie skończone fragmenty mają ten sam model. Dla różnych $T_0$ wolno wybierać inne struktury, a nawet struktury o różnych liczbach elementów. Wniosek twierdzenia o zwartości polega właśnie na możliwości znalezienia \emph{jednej} struktury dla całego $T$.
\end{uwaga}

\section{Twierdzenie o zwartości}\label{sec:compactness-main}
\begin{twierdzenie}[Zwartość logiki pierwszego rzędu]\label{thm:compactness}
Niech $\mathcal L$ będzie dowolnym językiem pierwszego rzędu z równością, a $T$ zbiorem zdań języka $\mathcal L$. Następujące warunki są równoważne:
\begin{enumerate}[label=(\roman*)]
\item $T$ ma model.
\item Każdy skończony podzbiór $T_0\subseteq T$ ma model.
\end{enumerate}
\end{twierdzenie}
\begin{proof}
Jeśli $T$ ma model $\mathfrak M$, to ten sam model spełnia każdy jego podzbiór; w szczególności każdy skończony. To dowodzi (i)$\Rightarrow$(ii).

Załóżmy (ii) i przypuśćmy, że $T$ nie ma modelu. Z Twierdzenia~\ref{thm:henkin-model-existence} (udowodnionego w poprzednim rozdziale) wnioskujemy kontrapozycyjnie, że $T$ jest syntaktycznie sprzeczny: istnieje formalna dedukcja $T\vdash_{\mathrm{ND}}\bot$. Na mocy Lematu~\ref{lem:henkin-finite-contradiction} dedukcja ta używa skończenie wielu zdań $T_0\subseteq T$, stąd $T_0\vdash_{\mathrm{ND}}\bot$. Twierdzenie o poprawności~\ref{thm:soundness} daje $T_0\models\bot$. Ponieważ $\bot$ jest fałszywe w każdej strukturze, żaden model nie spełnia $T_0$. Otrzymaliśmy sprzeczność z (ii). Zatem $T$ ma model.
\end{proof}

\begin{uwaga}[Skąd bierze się zwartość?]
Uzasadnienie używa \emph{dwóch} własności: pełności (brak modelu wykrywamy formalnym dowodem sprzeczności) oraz skończoności formalnego dowodu (w dowodzie może pojawić się tylko skończenie wiele przesłanek). Sama skończoność dowodów bez pełności nie wystarczyłaby do semantycznej zwartości. Rozróżnienie to było ważne już w rozdziale~\ref{ch:soundness-fol}.
\end{uwaga}

\begin{twierdzenie}[Skończone wsparcie konsekwencji semantycznej]\label{thm:compactness-finite-entailment}
Niech $\mathcal L$ będzie językiem pierwszego rzędu z równością, $\Gamma$ dowolnym zbiorem formuł tego języka (dopuszczamy zmienne wolne), a $\varphi$ formułą języka $\mathcal L$. Jeśli $\Gamma\models\varphi$, to istnieje skończony zbiór $\Gamma_0\subseteq\Gamma$ taki, że $\Gamma_0\models\varphi$.
\end{twierdzenie}
\begin{proof}
Twierdzenie~\ref{thm:godel-completeness} daje $\Gamma\vdash_{\mathrm{ND}}\varphi$. Dowód formalny jest skończonym drzewem, zatem na mocy Twierdzenia~\ref{thm:nd-finite-support} istnieje skończony $\Gamma_0\subseteq\Gamma$ spełniający $\Gamma_0\vdash_{\mathrm{ND}}\varphi$. Poprawność daje $\Gamma_0\models\varphi$.
\end{proof}

\begin{przyklad}[Dlaczego zwartość nie mówi o jednym skończonym modelu]
Zdania „istnieje przynajmniej $n$ różnych elementów”, osobno dla każdej liczby $n\ge1$, są jednocześnie spełnialne w $\N$. Jednak żaden \emph{skończony} model nie spełnia wszystkich tych zdań. Twierdzenie o zwartości obiecuje istnienie modelu, nie zaś istnienie modelu skończonego.
Kiedy próbujemy znaleźć jeden skończony model, zawsze możemy wybrać
zdanie wymagające co najmniej $n$ elementów, dla $n$ większego niż liczba elementów rozważanego modelu. Zdania te oznaczymy poniżej przez $\eta_n$.
Właśnie dlatego samo spełnienie każdego skończonego fragmentu
nie gwarantuje żadnego \emph{wspólnego modelu skończonego}.
\end{przyklad}

\section{Nieskończone modele z warunków skończonych}\label{sec:compactness-infinite}
\begin{definicja}[Zdania „co najmniej $n$ elementów”]\label{def:compactness-at-least-n}
Dla każdej liczby naturalnej $n\ge1$ zdefiniujmy zdanie $\eta_n$ następująco.
Dla $n=1$ kładziemy $\eta_1:=\exists x_1\,(x_1=x_1)$.
Dla $n\ge2$ piszemy
\[
\eta_n:=\exists x_1\ldots\exists x_n\;\bigwedge_{1\le i<j\le n}\neg(x_i=x_j).
\]
Symbol $\bigwedge$ oznacza koniunkcję wszystkich \emph{skończenie wielu}
formuł wskazanych przez indeksy: np.
$\eta_2=\exists x_1\exists x_2\,\neg(x_1=x_2)$,
a $\eta_3$ wymaga trzech parami różnych elementów.
Zapis $\bigwedge$ jest skrótem w metajęzyku, a nie nowym
spójnikiem języka pierwszego rzędu.
\end{definicja}

\begin{twierdzenie}[Istnienie modeli nieskończonych]\label{thm:compactness-infinite-model}
Niech $T$ będzie zbiorem zdań w języku pierwszego rzędu z równością. Jeżeli dla każdej liczby naturalnej $n\ge1$ istnieje model $T$ mający przynajmniej $n$ elementów, to $T$ ma model nieskończony.
\end{twierdzenie}
\begin{proof}
Rozważmy $T'=T\cup\{\eta_n:n\ge1\}$. Niech $F\subseteq T'$ będzie dowolnym skończonym zbiorem zdań. Wśród warunków $\eta_n$ należących do $F$ istnieje największy indeks $N$ (jeśli żaden taki warunek nie występuje, przyjmujemy $N=1$). Z założenia istnieje model $\mathfrak M_N$ zbioru $T$ o co najmniej $N$ elementach. Model ten spełnia wszystkie zdania z $F\cap T$ oraz wszystkie wymienione w $F$ zdania $\eta_n$, bo jeśli ma co najmniej $N$ elementów, to ma także co najmniej $n$ elementów dla $n\le N$. Zatem każdy skończony fragment $T'$ ma model. Zwartość daje model $\mathfrak M$ całego $T'$. Gdyby jego uniwersum miało dokładnie $m$ elementów, zdanie $\eta_{m+1}$ byłoby w nim fałszywe. Sprzeczność. Model $\mathfrak M$ jest więc nieskończony i spełnia $T$.
\end{proof}

\begin{wniosek}[Skończoności nie definiuje jedno zdanie pierwszego rzędu]\label{cor:compactness-no-finiteness-sentence}
Nie istnieje zdanie $\mathrm{Fin}$ języka zawierającego jedynie równość, które jest prawdziwe dokładnie w strukturach mających skończone, niepuste uniwersum.
\end{wniosek}
\begin{proof}
Załóżmy, że takie zdanie istnieje. Zbiór $\{\mathrm{Fin}\}\cup\{\eta_n:n\ge1\}$ jest skończenie spełnialny: jeśli wybierzemy tylko skończenie wiele warunków liczności, możemy wziąć dostatecznie duży skończony zbiór. Zwartość dałaby model całego zbioru zdań. Musiałby on jednocześnie być skończony (bo spełnia $\mathrm{Fin}$) i mieć co najmniej $n$ elementów dla każdego $n$ (bo spełnia wszystkie $\eta_n$). To niemożliwe.
\end{proof}

\begin{uwaga}[Czego ten wniosek nie mówi]
Dla ustalonego $n$ można w logice pierwszego rzędu opisać dokładnie $n$-elementowe zbiory, wymagając $n$ różnych elementów i tego, że każdy element jest jednym z nich. Nie istnieje natomiast \emph{jedno} zdanie, które byłoby prawdziwe dla wszystkich i tylko skończonych, niepustych struktur. Różnica polega na żądaniu jednej formuły działającej dla dowolnego rozmiaru.
\end{uwaga}

\section{Model arytmetyki z elementem większym od wszystkich standardowych liczb}\label{sec:compactness-nonstandard}
\begin{definicja}[Teoria pełna standardowych liczb naturalnych]\label{def:compactness-theory-n}
Niech $\mathcal L_{\mathrm{ar}}=\{0,S,+,\cdot,<\}$ oznacza język arytmetyki, w którym $0$ jest stałą, $S$ symbolem jednoargumentowej funkcji, $+$ i $\cdot$ symbolami funkcji dwuargumentowych, a $<$ symbolem relacji dwuargumentowej. Przez $\mathfrak N$ oznaczamy jego standardową strukturę o uniwersum $\N=\{0,1,2,\ldots\}$. Zbiór
\[
\operatorname{Th}(\mathfrak N)=\{\sigma:\sigma\text{ jest zdaniem języka }\mathcal L_{\mathrm{ar}}\text{ i }\mathfrak N\models\sigma\}
\]
nazywamy \emph{pełną teorią struktury $\mathfrak N$}. Przez $\overline n$ oznaczamy term $S^n(0)$ (np. $\overline0=0$, $\overline2=S(S(0))$).
\end{definicja}

\begin{twierdzenie}[Niestandardowy element w modelu teorii $\mathfrak N$]\label{thm:compactness-nonstandard}
Istnieje struktura $\mathfrak M$ języka $\mathcal L_{\mathrm{ar}}\cup\{c\}$,
w której $c$ jest nowym symbolem stałej, taka że
redukt $\mathfrak M\upharpoonright\mathcal L_{\mathrm{ar}}$
(zapomnienie interpretacji $c$) spełnia
$\operatorname{Th}(\mathfrak N)$ oraz dla każdego $n\in\N$
zachodzi $\mathfrak M\models(\overline n<c)$.
Redukt ten nie jest izomorficzny ze standardową strukturą $\mathfrak N$.
\end{twierdzenie}
\begin{proof}
Rozszerzmy język arytmetyki o nową stałą $c$ i rozważmy zbiór zdań
\[
T=\operatorname{Th}(\mathfrak N)\cup\{\overline n<c:n\in\N\}.
\]
Ustalmy dowolny skończony fragment $F\subseteq T$. Występuje w nim tylko skończenie wiele wymagań $\overline n<c$, powiedzmy dla indeksów $n_1,\ldots,n_k$. Wybierzmy zwykłą liczbę naturalną $m>\max\{n_1,\ldots,n_k\}$ (gdy żadnego wymagania nie ma, weźmy $m=0$). Interpretujmy $c$ w strukturze $\mathfrak N$ jako $m$. Wszystkie zdania z $F\cap\operatorname{Th}(\mathfrak N)$ są prawdziwe, ponieważ były prawdziwe w $\mathfrak N$ niezależnie od wartości nowej stałej. Spełnione są też wszystkie żądania $\overline n<c$ z $F$. Każdy skończony fragment $T$ ma więc model. Zwartość daje model $\mathfrak M$ całego $T$.

Gdyby redukt $\mathfrak M\upharpoonright\mathcal L_{\mathrm{ar}}$
był izomorficzny z $\mathfrak N$, to obrazem elementu
$c^{\mathfrak M}$ pod takim izomorfizmem byłaby pewna liczba naturalna $m$.
Izomorfizm zachowuje interpretacje $0$, $S$ i $<$,
a więc zdanie $\overline m<c$ w $\mathfrak M$ dawałoby
w $\mathfrak N$ nierówność $m<m$. To niemożliwe.
Zatem wskazane struktury nie są izomorficzne.
\end{proof}

\begin{uwaga}[Dlaczego nie powstaje sprzeczność z arytmetyką?]
W modelu $\mathfrak M$ interpretacja stałej $c$ jest większa od wartości każdego numerału $\overline n$ dla \emph{standardowego} $n$, a więc nie jest równa wartości żadnego z tych numerałów. Nie znaczy to, że model zawiera „liczbę większą od wszystkich swoich elementów”. Formuła $\forall x\,(x<c)$ w $\mathfrak M$ jest fałszywa, bo dla $x=c$ otrzymalibyśmy $c<c$. Zbiór warunków $\overline n<c$ zawiera osobne zdanie dla każdego \emph{zewnętrznie} ustalonego $n$, a nie jedno zdanie z kwantyfikatorem po wszystkich elementach $\mathfrak M$.
\end{uwaga}

\section{Kolorowanie grafu przez zwartość}\label{sec:compactness-coloring}
W poprzednich przykładach wymagania dotyczyły liczności. Pokażemy teraz zastosowanie kombinatoryczne, w którym każdy warunek odnosi się do skończonego fragmentu grafu, lecz wynik obejmuje graf dowolnie duży.

\begin{definicja}[Właściwe kolorowanie]\label{def:compactness-coloring}
Niech $G=(V,E)$ będzie grafem nieskierowanym bez pętli, czyli $E\subseteq V\times V$, $(v,w)\in E\iff(w,v)\in E$ oraz $(v,v)\notin E$, i niech $k\ge1$ będzie liczbą naturalną. \emph{Właściwym kolorowaniem $k$ kolorami} nazywamy funkcję $\kappa:V\to\{1,\ldots,k\}$ taką, że jeśli $(v,w)\in E$, to $\kappa(v)\ne\kappa(w)$.
\end{definicja}

\begin{twierdzenie}[Kolorowanie grafu na podstawie jego skończonych podgrafów]\label{thm:compactness-graph-coloring}
Niech $G=(V,E)$ będzie grafem nieskierowanym bez pętli, a $k\ge1$ ustaloną liczbą naturalną. Jeżeli każdy skończony podgraf $G$ ma właściwe kolorowanie $k$ kolorami, to cały graf $G$ również ma właściwe kolorowanie $k$ kolorami.
\end{twierdzenie}
\begin{proof}
Dla każdego wierzchołka $v\in V$ wprowadzamy nowy symbol stałej $c_v$, a dla każdego koloru $i\in\{1,\ldots,k\}$ nowy jednoargumentowy symbol relacyjny $P_i$. Dodajemy następujące zdania:
\begin{enumerate}[label=(\roman*)]
\item dla każdego $v\in V$ zdanie $P_1(c_v)\lor\cdots\lor P_k(c_v)$, wymagające co najmniej jednego koloru;
\item dla każdego $v$ i dla wszystkich $i<j$ zdanie $\neg(P_i(c_v)\land P_j(c_v))$, zabraniające dwóch różnych kolorów naraz;
\item dla każdej krawędzi $(v,w)\in E$ i każdego $i$ zdanie $\neg(P_i(c_v)\land P_i(c_w))$, zabraniające wspólnego koloru końcom krawędzi.
\end{enumerate}
Oznaczmy cały zbiór tych zdań przez $T_G$. Weźmy dowolny jego skończony podzbiór $F$. Pojawia się w nim tylko skończenie wiele nazw $c_v$; odpowiadające im wierzchołki tworzą skończony zbiór $W$. Z założenia podgraf indukowany przez $W$ można pokolorować $k$ kolorami. Wystarczy zatem wziąć strukturę o uniwersum $\{1,\ldots,k\}$, każdą stałą $c_v$ wymienioną w $F$ zinterpretować jako kolor wierzchołka $v$, pozostałe stałe zinterpretować dowolnie, a $P_i$ interpretować jako jednoelementowy zbiór $\{i\}$. Wtedy wszystkie zdania z $F$ są spełnione. Zwartość daje model $\mathfrak M$ całego $T_G$.

Dla każdego $v\in V$ z (i) i (ii) wynika istnienie dokładnie jednego indeksu $i$ z $\mathfrak M\models P_i(c_v)$. Definiujemy $\kappa(v)=i$. Jeżeli $(v,w)\in E$ i $\kappa(v)=\kappa(w)=i$, mielibyśmy w $\mathfrak M$ prawdziwe równocześnie $P_i(c_v)$ i $P_i(c_w)$, sprzecznie z (iii). Otrzymaliśmy więc właściwe kolorowanie całego $G$.
\end{proof}

\begin{przyklad}[Kolorowanie nieskończonej ścieżki]
Dla grafu o wierzchołkach $\Z$ i krawędziach między $n$ i $n+1$ każde skończone pasmo ma kolorowanie dwoma kolorami. Zwartość gwarantuje kolorowanie całej ścieżki. W tym prostym przypadku można je też zapisać bezpośrednio, kładąc $\kappa(n)=1$ dla parzystych $n$ i $\kappa(n)=2$ dla nieparzystych. Znaczenie twierdzenia nie polega więc na trudności tego konkretnego grafu, lecz na tym, że działa ono dla dowolnego, nawet nieprzeliczalnego zbioru wierzchołków.
\end{przyklad}

\section{Własność, której nie można opisać zdaniami pierwszego rzędu}\label{sec:compactness-connectivity}
\begin{definicja}[Spacer, ścieżka i spójność grafu]\label{def:compactness-graph-connectivity}
Niech $G=(V,E)$ będzie grafem nieskierowanym bez pętli, niech $v,w\in V$, a $m\in\N$. \emph{Spacerem długości $m$} od $v$ do $w$ nazywamy ciąg wierzchołków $v_0,\ldots,v_m$ taki, że $v_0=v$, $v_m=w$ oraz $(v_{j-1},v_j)\in E$ dla $j=1,\ldots,m$. Wierzchołki mogą się powtarzać; dla $v=w$ dopuszczamy spacer długości $0$. Graf nazywamy \emph{spójnym}, jeżeli każde dwa wierzchołki można połączyć spacerem skończonej długości. Jest to równoważne istnieniu ścieżki bez powtórzeń wierzchołków: powtarzające się odcinki spaceru można kolejno usuwać.
\end{definicja}

\begin{twierdzenie}[Spójność grafów nie jest aksjomatyzowalna w logice pierwszego rzędu]\label{thm:compactness-not-connectivity}
Nie istnieje zbiór $\Sigma$ zdań języka pierwszego rzędu $\mathcal L=\{E\}$ z równością, gdzie $E$ jest dwuargumentowym symbolem relacyjnym, którego modelami byłyby dokładnie wszystkie niepuste, nieskierowane grafy bez pętli, które są spójne.
\end{twierdzenie}
\begin{proof}
Przypuśćmy, że taki zbiór zdań $\Sigma$ istnieje. Dodajmy dwie nowe stałe $a,b$. Dla każdego $n\ge0$ można zapisać zdanie $\rho_n(a,b)$ mówiące, że \emph{nie} ma spaceru długości co najwyżej $n$ z $a$ do $b$. Symbole $\mathrm{Walk}_m$ oznaczają tu formuły opisujące spacery, nie ścieżki bez powtórzeń wierzchołków. Dokładniej, dla $m=0$ kładziemy
$\mathrm{Walk}_0(a,b):=(a=b)$, a dla $m=1$:
$\mathrm{Walk}_1(a,b):=E(a,b)$. Dla $m\ge2$ przyjmijmy $z_0:=a$ i $z_m:=b$ oraz połóżmy
\[
\mathrm{Walk}_m(a,b):=\exists z_1\ldots\exists z_{m-1}\;
\bigwedge_{j=1}^{m}E(z_{j-1},z_j).
\]
Koniunkcja jest skończona dla każdego ustalonego $m$ i opisuje spacer o $m$ krawędziach. Na przykład $\mathrm{Walk}_2(a,b)$ to $\exists z_1\,(E(a,z_1)\land E(z_1,b))$.
Wreszcie $\rho_n(a,b)$ jest zdaniem
$\neg(\mathrm{Walk}_0(a,b)\lor\cdots\lor\mathrm{Walk}_n(a,b))$.
Jest to skończona alternatywa dla każdego ustalonego $n$;
nie piszemy tutaj nieskończonej formuły.

Rozważmy teorię $T=\Sigma\cup\{\rho_n(a,b):n\ge0\}$. Dowolny skończony fragment $F\subseteq T$ zawiera tylko skończenie wiele wymagań $\rho_n$, a zatem, jeśli takie wymagania występują, istnieje największy ich indeks $N$; w przeciwnym przypadku przyjmijmy $N=0$. Wybierzmy skończony graf będący zwykłą ścieżką o $N+2$ krawędziach, a jej końce interpretujmy jako $a$ i $b$. Jest to graf spójny, więc spełnia wszystkie zdania z $\Sigma$, oraz nie ma spaceru między końcami długości co najwyżej $N$, więc spełnia również $F$. Zwartość daje model $\mathfrak M$ całego $T$. Ponieważ $\mathfrak M\models\Sigma$, graf $\mathfrak M$ musi być spójny. Istnieje zatem spacer długości pewnego $m$ od $a^{\mathfrak M}$ do $b^{\mathfrak M}$. Ale $\mathfrak M\models\rho_m(a,b)$ przeczy istnieniu takiego spaceru. To sprzeczność.
\end{proof}

\begin{uwaga}[Zależność od języka]
Twierdzenie nie mówi, że spójności nie można \emph{zdefiniować w ogóle}. Można ją opisać w silniejszych językach, np. dopuszczających odpowiednie kwantyfikowanie po podzbiorach wierzchołków, albo za pomocą nieskończenie długiej alternatywy „istnieje ścieżka długości $0$ lub $1$ lub $2$ albo $\ldots$”. Takie wyrażenie nie jest jednak jedną formułą zwykłej logiki pierwszego rzędu: każda jej formuła jest skończonym napisem.
\end{uwaga}


\begin{twierdzenie}[Nie istnieje skończona aksjomatyzacja nieskończoności]\label{thm:compactness-infinity-not-finitely-axiomatizable-47}
W języku zawierającym wyłącznie równość nie istnieje skończony zbiór
zdań $F$, którego modelami byłyby dokładnie wszystkie struktury
o nieskończonym, niepustym uniwersum.
\end{twierdzenie}
\begin{proof}
Przypuśćmy, że taki skończony zbiór $F$ istnieje. Niech $\theta$
będzie koniunkcją wszystkich zdań z $F$; jeśli $F=\varnothing$,
niech $\theta$ będzie zdaniem $\forall x\,(x=x)$.
Wówczas $\theta$ jest prawdziwe dokładnie w strukturach nieskończonych.
Dla każdego $n\ge1$ rozważmy teorię
$T=\{\neg\theta\}\cup\{\eta_n:n\ge1\}$, gdzie $\eta_n$
zostały zdefiniowane w Definicji~\ref{def:compactness-at-least-n}.
Dowolny skończony fragment $T$ zawiera tylko skończenie wiele zdań
$\eta_n$, zatem spełnia go odpowiednio duża struktura skończona:
w takiej strukturze $\neg\theta$ jest prawdziwe na mocy założenia.
Zwartość daje model całego $T$. Ponieważ spełnia on wszystkie
$\eta_n$, jego uniwersum jest nieskończone, więc spełnia $\theta$.
Jednocześnie spełnia $\neg\theta$, co jest sprzecznością.
\end{proof}

\begin{przyklad}[Teoria nieskończoności a pojedyncze zdanie]\label{ex:compactness-infinity-theory-47}
Zbiór $I=\{\eta_n:n\ge1\}$ jest nieskończoną teorią pierwszego
rzędu. Struktura spełnia $I$ dokładnie wtedy, gdy ma nieskończenie
wiele elementów: jeśli jest nieskończona, dla każdego ustalonego
$n$ można wybrać $n$ różnych elementów; jeśli ma $m$ elementów,
nie spełnia $\eta_{m+1}$. Poprzednie twierdzenie mówi, że nie
można zastąpić całego $I$ żadnym skończonym zbiorem zdań,
choć każdy jego pojedynczy warunek jest formułą pierwszego rzędu.
Nie należy więc mylić aksjomatyzowalności przez teorię
z definiowalnością przez jedno zdanie.
\end{przyklad}

\section{Granice twierdzenia i dalsza droga}
Zwartość dotyczy \emph{spełnialności} zdań pierwszego rzędu; nie podaje algorytmu budowania modeli ani rozstrzygania teorii. Pełność logiki pierwszego rzędu dotyczy konsekwencji prawdziwych we \emph{wszystkich modelach} przesłanek, nie dowodliwości każdego zdania prawdziwego w standardowej arytmetyce. Twierdzenia Gödla o niezupełności dotyczą zdań nierozstrzygniętych przez konkretne efektywnie aksjomatyzowane teorie. Następne rozdziały rozwiną analizę modeli i obliczalności.